% DRAFT ONLY: independent root audit pending. 2026-11-19 17:41 UTC.
% A.126 is reserved for the closed simply-connected dimension-five case.
\begin{saresult}{A.127}{Proved under stated hypotheses}{The projective plane has trivial concordance component group}
For the standard PL projective plane and every definable model over
any real closed field,
\[
 \pi_0P_{\rm PL}(\mathbb {RP}^2)
 =\pi_0P_{\rm def}(\mathbb {RP}^2)=0.
\]
The bottom is fixed pointwise and the top is free. Together with
A.123 and A.125 this computes the component group for every closed
surface. No higher parametric comparison is asserted.
\end{saresult}
\begin{proof}
Put $M=\mathbb {RP}^2\times I$. We work over the real field first.
Every homeomorphism and every isotopy constructed below is jointly
PL, including the twists used to change the annulus winding.

We use two elementary consequences of the three-dimensional
PL Schoenflies theorem, whose original sphere statement is
\cite[p.~6, with proof pp.~7--8]{Alexander1924}.
First, $M$ is irreducible. Its orientation double cover is
$S^2\times I$, with deck involution antipodal on $S^2$ and identity
on $I$. An embedded PL sphere in $\operatorname{int}M$ lifts to two
disjoint spheres exchanged by this involution. Cap the two ends of
$S^2\times I$ by balls to obtain $S^3$. For each lifted sphere,
Schoenflies gives the following dichotomy: either the two end caps
lie on the same side, and the other side is a ball lying wholly in
the open cylinder, or the sphere separates the two ends.
Two disjoint spheres separating the ends are strictly ordered from
bottom to top. An involution preserving each end cannot exchange
them. Hence the lifts bound interior balls. Balls with disjoint
boundaries in this cylinder are disjoint or nested; strict nesting
of the two balls is impossible because the involution has order two.
The balls are therefore disjoint, and either projects to a ball
bounded by the original sphere. This proves irreducibility.

Second, two proper PL disks in a PL three-ball with the same
boundary curve are ambient isotopic relative to the entire boundary
of the ball. To see this without an intersection-minimization
assumption, each disk separates the ball into two PL balls, by
Schoenflies applied to the disk joined to either boundary hemisphere.
Choose a PL homeomorphism between the disks, fixed on their common
boundary. On each complementary ball, glue this map to the identity
on the corresponding boundary hemisphere. The resulting boundary
sphere homeomorphism extends across that ball by coning. The two
extensions glue to a homeomorphism of the original ball fixed on
its entire boundary and taking the first disk to the second. The
Alexander isotopy contracts this ball homeomorphism relative to its
boundary, proving the assertion. The same conical formula shows
that this isotopy is jointly PL, as in A.123.

We describe the standard annuli and suitable PL twists explicitly.
Let $D$ be a centrally symmetric polygonal disk, identify its
boundary with $\mathbb R/2\mathbb Z$ by an edgewise linear
coordinate $s$, and make the identification compatible with
$x\mapsto-x$ by $s\mapsto s+1$. The quotient of $D$ by antipodal
identification on its boundary is $\mathbb {RP}^2$.
The image of $\partial D$ is a one-sided projective line $c$.
The image of a diameter of $D$ joining the points $s=0$ and $s=1$
is another one-sided projective line $a$, meeting $c$ once. Put
\[
 C=c\times I,\qquad A_0=a\times I.
\]
Both are proper one-sided annuli. Cutting $M$ along either one
gives $D^2\times I$, a three-ball. In particular cutting along $C$
is exactly the quotient description $D\times I\to M$.

We first construct a jointly PL disk isotopy from the identity to
the exact linear antipodal map $J:x\mapsto-x$, such that every
boundary map commutes with $J$. The boundary translation
$c_u(s)=s+u$ on $\mathbb R/2\mathbb Z$ is jointly PL. On a PL
collar $\partial D\times[0,\delta]$, with $v=0$ the outer boundary,
extend it by
\[
 (s,v)\longmapsto
 \bigl(s+\max\{0,u-v/\delta\},v\bigr),
 \qquad 0\leq u\leq1,
\]
and use the identity outside the collar. This gives a jointly PL
isotopy $r_u$ of $D$ whose final boundary map is $J$. The map
$b=J^{-1}r_1$ fixes $\partial D$ pointwise. By the relative disk
Alexander isotopy, join $b$ to the identity through
boundary-fixed PL maps $b_u$. Concatenating $r_u$ with $Jb_u$
gives a jointly PL isotopy from the identity to exactly $J$.
On the first half its boundary map is translation by a number from
zero to one; on the second half it is constantly translation by
one. In particular every boundary map commutes with $s\mapsto s+1$.

Repeat this path, composing successive copies with $J^j$, to obtain
for each positive integer $k$ a jointly PL path $q_t$ from the
identity to $J^k$. For negative $k$ use the inverse path on each
copy; for $k=0$ use the constant path. The lift of its boundary-translation amount starts at zero
and ends at $k$, and is constant only on the Alexander-correction
parts. Since all boundary maps commute with the antipodal
identification, $q_t$ descends to a jointly PL isotopy $\bar q_t$
of $\mathbb {RP}^2$. Define
\[
 T_k(x,t)=(\bar q_t(x),t),\qquad
 T_{k,u}(x,t)=(\bar q_{\min\{u,t\}}(x),t).             \tag{A.127.1}
\]
The second family is a jointly PL bottom-fixed isotopy from the
identity to $T_k$. Set $A_k=T_k(A_0)$. Its intersection with $C$
is the graph of the boundary-translation amount modulo one, hence
a proper arc of winding number $k$ with the same endpoints as
$A_0\cap C$. Its bottom and top boundary curves are both $a$,
since $q_0$ is the identity and $q_1=J^k$ preserves the diameter.
This supplies actual PL twists and their PL contractions; no polar
rotation formula is being called PL.

Let $h\in P_{\rm PL}(\mathbb {RP}^2)$. Its upper restriction is
homotopic to the identity through
$x\mapsto\operatorname{pr}_1h(x,t)$. The surface homotopy-isotopy
theorem \cite[Theorem~6.3, p.~99]{Epstein1966} gives a PL isotopy
$a_u$ of the upper map $a$ to the identity. Put $c_u=a_ua^{-1}$
and extend this inverse correction by
$C_u(x,t)=(c_{\min\{u,t\}}(x),t)$. This is a jointly PL
bottom-fixed isotopy starting at the identity, and $C_1h$ has
upper restriction $a^{-1}a=\operatorname{id}$.
After this bottom-fixed PL isotopy we may assume that
$h$ fixes both boundary components pointwise.

Put $A=h(A_0)$. Arrange $A$ and $C$ in PL general position
relative to the boundary. Their boundaries meet in exactly two
points, one on each end. Thus their intersection consists of one
proper spanning arc and finitely many circles. Every such circle
is inessential in each annulus: an essential circle in an annulus
must meet every arc joining its two boundary components, whereas
the intersection components are disjoint. Choose a circle innermost
on $C$. It bounds a disk $D_C\subset C$ with interior disjoint
from $A$, and also a disk $D_A\subset A$. Their union is an
embedded PL sphere in the interior of $M$, and bounds a ball $B$
by irreducibility. The surface
$(A\setminus\operatorname{int}D_A)\setminus\partial D_A$ is
connected, contains the two original boundary circles of $A$ on
$\partial M$, and is disjoint from this sphere. It consequently
lies outside $B$. In particular the
interior of $B$ misses $A$. Pushing $D_A$ across this ball and
slightly past $D_C$ removes the selected circle and any circles
inside $D_A$, introduces no new intersection components, and is
supported away from $\partial M$. Repeating finitely often leaves
a single spanning arc.

This remaining arc has an integer winding number $k$ in $C$,
relative to its fixed endpoints. The relative arc theorem
\cite[Theorem~3.1, p.~89]{Epstein1966} isotopes it to
$A_k\cap C$, keeping $\partial C$ fixed. Extend this isotopy to
the ambient manifold, relative to $\partial M$, using the normal
interval-bundle construction described below. After this change,
cut along $C$. The annuli $A$ and $A_k$ each cut to a proper PL
disk in the ball $D\times I$, and their boundary curves agree:
their two side arcs are the two lifts of $A_k\cap C$, and their
end arcs are the common diameter on the top and bottom.
The disk assertion already proved gives an ambient PL isotopy
fixed on the entire cut-ball boundary taking one disk to the other.
It descends to $M$. Therefore an ambient PL isotopy relative to
$\partial M$ has changed $h$ to a map $h_1$ with
$h_1(A_0)=A_k$.

Compose with $T_k^{-1}$ and put $g=T_k^{-1}h_1$. This is still PL,
fixes the bottom, and preserves $A_0$ setwise. It is connected to
$h$ by a bottom-fixed PL isotopy, by the explicit path for $T_k$.
Its restriction $f=g|_{A_0}$ is a bottom-fixed annulus
homeomorphism. The circle calculation in A.123 gives a jointly PL
isotopy $f_u$ from $f$ to the identity, fixed on the bottom circle.
Set $c_u=f_uf^{-1}$, an isotopy from the identity to $f^{-1}$.

Here is the promised extension, including the one-sided case.
A boundary-compatible normal neighborhood of the proper annulus
$A_0$ is
\[
 N=(\widetilde A_0\times[-\delta,\delta])/
       ((p,v)\sim(dp,-v)),
\]
where $\widetilde A_0$ is its normal-orientation double cover and
$d$ is the deck involution. Lift $c_u$ starting at the identity to
$\widetilde c_u$. The lifts commute with $d$ and fix every lift
of a boundary point fixed by $c_u$. Define
\[
 B_u([p,v])=
 \left[\widetilde c_{\tau(u,|v|)}(p),v\right],\qquad
 \tau(u,r)=\max\{0,u-r/\delta\}.                    \tag{A.127.2}
\]
This is well defined, is the identity for $|v|=\delta$, and has
inverse given by the same formula with the inverse lifted family.
It therefore extends by the identity outside $N$ to an ambient
isotopy. Covering lifts and composition of the finite PL maps show
that it is jointly PL. It fixes the whole bottom because $c_u$
fixes its bottom circle. For the preceding arc isotopy in $C$ the
same construction fixes both boundary components, as that isotopy
fixes both boundary circles. Thus no two-sided normal framing was
assumed at either use.

The map $g'=B_1g$ now fixes $A_0$ pointwise. Its action on the two
normal directions is locally constant along the connected annulus.
It is the identity at the bottom, since the entire bottom surface
is fixed. Hence it preserves the normal directions everywhere.
Cutting along $A_0$ consequently lifts $g'$ to a PL homeomorphism
of $D^2\times I$ fixing its bottom and its whole side pointwise.
The ball calculation of A.123 gives an isotopy to the identity
relative to these sets: first cancel the free top-disk map by its
boundary-relative Alexander isotopy and a collar extension, and
then Alexander-contract the resulting full-boundary-fixed
three-ball map. This isotopy descends to $M$. We have proved that
the original PL map $h$ is jointly PL isotopic to the identity through
bottom-fixed maps. This proves the real PL assertion. A.115 and the
relative component argument of A.113 give the assertion over every
real closed field and for every definable model.

The integer $k$ was used only to select an arc model after a chosen
intersection simplification. We neither asserted that it is an
ambient-isotopy invariant nor that $T_1$ fixes the upper projective
plane. Those stronger assertions are unnecessary and false in the
literal rotation model.
For the precise source correction, the displayed twist in
\cite[Appendix~A, Theorem~A.6, pp.~614--615]{BoyleChen2024}
sends the upper interior point $(1/2,0)$ of the disk model to
$(-1/2,0)$. They are distinct quotient points, because only
boundary points are identified. Thus that literal twist does not
belong to the full-boundary-fixed group claimed there. The exact
two-point check is recorded with the source audit; the assertion of
that theorem is not an input to the present free-top proof.
\end{proof}
